\section{快速上手与基础操作}
%% ============================================================

掌握以下基础操作可以显著提升日常使用效率。

\subsection{设置 cc 别名（Tip \#1）}

在 \texttt{\textasciitilde/.zshrc}（或 \texttt{\textasciitilde/.bashrc}）中添加别名，用 \texttt{cc} 代替 \texttt{claude} 启动会话，同时跳过所有权限确认弹窗。

\begin{lstlisting}[caption={cc 别名配置}, language=bash]
alias cc="claude --dangerously-skip-permissions"
\end{lstlisting}

\begin{warningbox}{谨慎使用 --dangerously-skip-permissions}
该 flag 的命名故意设计得"吓人"。只有在你完全理解 Claude Code 能对你的代码库做什么之后，才应启用此选项。它会跳过所有操作确认，包括文件删除和 shell 命令执行。
\end{warningbox}

\subsection{用 ! 前缀直接执行 Bash 命令（Tip \#2）}

在 Claude Code 交互界面中输入 \texttt{!git status} 或 \texttt{!npm test}，命令会立即执行，输出会进入上下文供 Claude 参考。比让 Claude 帮你跑命令更快。

\subsection{Esc 中断与 Esc+Esc 回滚（Tip \#3）}

\begin{itemize}
  \item \textbf{Esc}：立即中断 Claude 的当前操作，不丢失上下文，可以马上重新引导。
  \item \textbf{Esc+Esc}（或 \texttt{/rewind}）：打开 checkpoint 列表，可恢复代码、对话或两者。四种恢复选项：代码+对话、仅对话、仅代码、从某个 checkpoint 开始摘要。
\end{itemize}

\begin{importantbox}{Checkpoint 的局限}
Checkpoint 只追踪文件编辑操作。通过 Bash 命令造成的变更（数据库迁移、文件系统操作等）不会被 checkpoint 捕获，无法回滚。
\end{importantbox}

恢复上次会话：\texttt{claude --continue} 恢复最近一次对话；\texttt{claude --resume} 打开会话选择器。

\subsection{常用快捷键汇总（Tips \#16, \#17, \#27）}

\begin{table}[H]
\centering
\begin{tabular}{lp{9cm}}
\toprule
\textbf{快捷键} & \textbf{功能} \\
\midrule
\texttt{Ctrl+S} & 暂存当前正在编写的 prompt，处理完其他事情后自动恢复 \\
\texttt{Ctrl+B} & 将长时间运行的 Bash 命令放入后台，Claude 继续工作 \\
\texttt{Ctrl+G} & 在文本编辑器中直接编辑 Claude 提出的计划 \\
\texttt{Shift+Tab} & 在 Normal / Auto-Accept / Plan 三种权限模式间切换 \\
\texttt{Esc} & 中断当前操作 \\
\texttt{Esc+Esc} & 打开 checkpoint 回滚菜单 \\
\bottomrule
\end{tabular}
\caption{Claude Code 常用快捷键}
\end{table}

\subsection{命名与着色会话（Tip \#47）}

运行多个并行会话时，用 \texttt{/rename auth-refactor} 给会话加标签，用 \texttt{/color red} 设置提示栏颜色（支持 red, blue, green, yellow, purple, orange, pink, cyan）。五秒钟的操作能避免你在错误的终端里输入命令。

\subsection{自定义 Spinner 动词（Tip \#50）}

Claude 思考时终端显示的 spinner 动词（如"Flibbertigibbeting..."）可以自定义。只需告诉 Claude："Replace my spinner verbs with Harry Potter spells."，Claude 会自动生成列表。虽然是小事，但能让等待变得有趣。

\subsection{本章小结}

基础操作的核心在于：用别名简化启动、用 \texttt{!} 前缀快速执行命令、用 Esc 系列快捷键控制流程、用命名和颜色管理多会话。这些习惯一旦养成，日常交互效率会有质的提升。

%% ============================================================
